\documentclass[10pt,a4paper]{amsart}
\usepackage{amsmath,amssymb,amsthm,graphicx}
\usepackage{verbatim}
\newenvironment{reference}{\comment}{\endcomment}
%\newenvironment{reference}{}{}
\newenvironment{slogan}{\comment}{\endcomment}
\newenvironment{history}{\comment}{\endcomment}

% For commutative diagrams we use Xy-pic
\usepackage[all]{xy}

% We use 2cell for 2-commutative diagrams.
\xyoption{2cell}
% Theorem environments.
%
\theoremstyle{plain}
\newtheorem{theorem}[subsection]{Theorem}
\newtheorem{proposition}[subsection]{Proposition}
\newtheorem{lemma}[subsection]{Lemma}

\theoremstyle{definition}
\newtheorem{definition}[subsection]{Definition}
\newtheorem{example}[subsection]{Example}
\newtheorem{exercise}[subsection]{Exercise}
\newtheorem{situation}[subsection]{Situation}

\theoremstyle{remark}
\newtheorem{remark}[subsection]{Remark}
\newtheorem{remarks}[subsection]{Remarks}

\numberwithin{equation}{subsection}

% Macros
%
\def\lim{\mathop{\mathrm{lim}}\nolimits}
\def\colim{\mathop{\mathrm{colim}}\nolimits}
\def\Spec{\mathop{\mathrm{Spec}}}
\def\Hom{\mathop{\mathrm{Hom}}\nolimits}
\def\Ext{\mathop{\mathrm{Ext}}\nolimits}
\def\SheafHom{\mathop{\mathcal{H}\!\mathit{om}}\nolimits}
\def\SheafExt{\mathop{\mathcal{E}\!\mathit{xt}}\nolimits}
\def\Sch{\mathit{Sch}}
\def\Mor{\mathop{\mathrm{Mor}}\nolimits}
\def\Ob{\mathop{\mathrm{Ob}}\nolimits}
\def\Sh{\mathop{\mathit{Sh}}\nolimits}
\def\NL{\mathop{N\!L}\nolimits}
\def\CH{\mathop{\mathrm{CH}}\nolimits}
\def\proetale{{pro\text{-}\acute{e}tale}}
\def\etale{{\acute{e}tale}}
\def\QCoh{\mathit{QCoh}}
\def\Ker{\mathop{\mathrm{Ker}}}
\def\Im{\mathop{\mathrm{Im}}}
\def\Coker{\mathop{\mathrm{Coker}}}
\def\Coim{\mathop{\mathrm{Coim}}}

% Boxtimes
%
\DeclareMathSymbol{\boxtimes}{\mathbin}{AMSa}{"02}

%
% Macros for moduli stacks/spaces
%
\def\QCohstack{\mathcal{QC}\!\mathit{oh}}
\def\Cohstack{\mathcal{C}\!\mathit{oh}}
\def\Spacesstack{\mathcal{S}\!\mathit{paces}}
\def\Quotfunctor{\mathrm{Quot}}
\def\Hilbfunctor{\mathrm{Hilb}}
\def\Curvesstack{\mathcal{C}\!\mathit{urves}}
\def\Polarizedstack{\mathcal{P}\!\mathit{olarized}}
\def\Complexesstack{\mathcal{C}\!\mathit{omplexes}}
% \Pic is the operator that assigns to X its picard group, usage \Pic(X)
% \Picardstack_{X/B} denotes the Picard stack of X over B
% \Picardfunctor_{X/B} denotes the Picard functor of X over B
\def\Pic{\mathop{\mathrm{Pic}}\nolimits}
\def\Picardstack{\mathcal{P}\!\mathit{ic}}
\def\Picardfunctor{\mathrm{Pic}}
\def\Deformationcategory{\mathcal{D}\!\mathit{ef}}
\usepackage{cleveref}
\crefname{originalreference}{original reference}{original references}
\Crefname{originalreference}{Original reference}{Original references}
\makeatletter
\newlabel{lemma-base-change-rig-surjective}{{0AQS}{source}{}{sourcepointer-1}{}}
\newlabel{lemma-base-change-rig-surjective@cref}{{[originalreference][0][]0AQS}{[1][0][]source}{}{}{}}
\makeatother
\crefname{originalreference}{original reference}{original references}
\Crefname{originalreference}{Original reference}{Original references}
\makeatletter
\newlabel{lemma-composition-rig-surjective}{{0AQR}{source}{}{sourcepointer-2}{}}
\newlabel{lemma-composition-rig-surjective@cref}{{[originalreference][0][]0AQR}{[1][0][]source}{}{}{}}
\makeatother
\crefname{originalreference}{original reference}{original references}
\Crefname{originalreference}{Original reference}{Original references}
\makeatletter
\newlabel{lemma-etale-covering-rig-surjective}{{0AQU}{source}{}{sourcepointer-3}{}}
\newlabel{lemma-etale-covering-rig-surjective@cref}{{[originalreference][0][]0AQU}{[1][0][]source}{}{}{}}
\makeatother
\crefname{originalreference}{original reference}{original references}
\Crefname{originalreference}{Original reference}{Original references}
\makeatletter
\newlabel{lemma-permanence-rig-surjective}{{0AQT}{source}{}{sourcepointer-4}{}}
\newlabel{lemma-permanence-rig-surjective@cref}{{[originalreference][0][]0AQT}{[1][0][]source}{}{}{}}
\makeatother
\title{Rig-surjectivity descends along etale covers}
\author{The Stacks Project Authors}
\date{stacks-project:restricted.tex:local-snapshot-20261012; GNU FDL 1.2 or later}
\begin{document}
\maketitle
\noindent\textit{Editorial note.} All delivered mathematical text is the author's own, copied from the pinned Stacks Project chapter file restricted.tex; no statement or proof step was rewritten, reordered or removed.\par\medskip
\noindent\textit{Editorial note.} The counted proof is the author's complete argument: the proof assembles base-change stability, etale covers being rig-surjective, composition and the factorization permanence lemma into the descent statement\par\medskip
\noindent\textit{Editorial note.} The label lemma-upshot names the counted statement, and the proof environment immediately following it in the source is delivered verbatim as its proof.\par\medskip
\noindent\textit{Editorial note.} The counted statement is carried by the author's own lemma environment in the chapter source.\par\medskip
\noindent\textit{Allowed standard prerequisites (editorial).}\begin{itemize}
\item base change of rig-surjective morphisms (lemma-base-change-rig-surjective of the same chapter)
\item etale covers are rig-surjective (lemma-etale-covering-rig-surjective of the same chapter) and rig-surjectivity composes (lemma-composition-rig-surjective of the same chapter)
\item permanence of rig-surjectivity along factorizations (lemma-permanence-rig-surjective of the same chapter)
\end{itemize}\medskip
\ifdefined\hypertarget\hypertarget{sourcepointer-1}{}\fi
\noindent\textit{Source reference.} Original reference 0AQS: sources/stacks-restricted/pinned-original-source.tex, source lines 5324--5324. Prerequisite pointer invoked at proof line(s) 5406; it names a standard fact of the same chapter that the counted argument uses, and every step of the counted proof built on it is retained verbatim in the delivered body.\par\medskip
\ifdefined\hypertarget\hypertarget{sourcepointer-2}{}\fi
\noindent\textit{Source reference.} Original reference 0AQR: sources/stacks-restricted/pinned-original-source.tex, source lines 5307--5307. Prerequisite pointer invoked at proof line(s) 5412; it names a standard fact of the same chapter that the counted argument uses, and every step of the counted proof built on it is retained verbatim in the delivered body.\par\medskip
\ifdefined\hypertarget\hypertarget{sourcepointer-3}{}\fi
\noindent\textit{Source reference.} Original reference 0AQU: sources/stacks-restricted/pinned-original-source.tex, source lines 5363--5363. Prerequisite pointer invoked at proof line(s) 5409; it names a standard fact of the same chapter that the counted argument uses, and every step of the counted proof built on it is retained verbatim in the delivered body.\par\medskip
\ifdefined\hypertarget\hypertarget{sourcepointer-4}{}\fi
\noindent\textit{Source reference.} Original reference 0AQT: sources/stacks-restricted/pinned-original-source.tex, source lines 5351--5351. Prerequisite pointer invoked at proof line(s) 5414; it names a standard fact of the same chapter that the counted argument uses, and every step of the counted proof built on it is retained verbatim in the delivered body.\par\medskip
\begin{lemma}
\label{lemma-upshot}
Let $S$ be a scheme. Let $f : X \to Y$ be a morphism of locally
Noetherian formal algebraic spaces which is locally of finite type.
Let $\{g_i : Y_i \to Y\}$ be a family of morphisms of formal
algebraic spaces which are representable by algebraic spaces and
\'etale such that $\coprod g_i$ is surjective.
Then $f$ is rig-surjective if and only if each
$f_i : X \times_Y Y_i \to Y_i$ is rig-surjective.
\end{lemma}

\begin{proof}
Namely, if $f$ is rig-surjective, so is any base change
(Lemma \ref{lemma-base-change-rig-surjective}).
Conversely, if all $f_i$ are rig-surjective, so is
$\coprod f_i : \coprod X \times_Y Y_i \to \coprod Y_i$.
By Lemma \ref{lemma-etale-covering-rig-surjective}
the morphism $\coprod g_i : \coprod Y_i \to Y$ is rig-surjective.
Hence $\coprod X \times_Y Y_i \to Y$ is rig-surjective
(Lemma \ref{lemma-composition-rig-surjective}).
Since this morphism factors through $X \to Y$ we see that $X \to Y$
is rig-surjective by Lemma \ref{lemma-permanence-rig-surjective}.
\end{proof}

\end{document}
